
\subsection{AlphaFold2 acceleration}
\label{sec:rw-alphafold}

AlphaFold2 inference optimization has been studied almost entirely on
NVIDIA GPUs. ParaFold separates the CPU-bound MSA search stage from
GPU inference and avoids redundant JAX recompilation across a batch of
proteins by sorting them by length, reporting an average 13.8$\times$
speedup over baseline AlphaFold2 \citep{zhong2022parafold}. FastFold
introduces Dynamic Axial Parallelism, a model-parallel scheme specific
to AlphaFold2's Evoformer, together with an activation-memory reduction
technique (AutoChunk); its arXiv version reports scaling to 512 GPUs
\citep{cheng2022fastfoldarxiv}. We do not place our 86.5\% parallel
efficiency on 8 chips (Section~\ref{sec:parallelism}) next to the
efficiency reported there: that figure is a \emph{training} scaling
result and ours is inference throughput, and the two are not directly
comparable \citep{cheng2024fastfold,cheng2022fastfoldarxiv}. APACE distributes
AlphaFold2's ensemble across hundreds of A100 GPUs to reduce
time-to-solution on real structures from weeks to minutes
\citep{apace2024}. ScaleFold and HelixFold scale AlphaFold2
\emph{training}, rather than inference, to large GPU clusters
\citep{scalefold2024,helixfold2022}. Uni-Fold reimplements AlphaFold2 in
PyTorch specifically because the original was TPU-trained and TPU
access is scarce, an explicit statement of the accessibility gap this
paper's TPU results speak to \citep{unifold}. OpenFold is a widely used GPU-targeted
reimplementation and ColabFold a widely used inference workflow around
the released models; ColabFold's roughly
90$\times$ batch speedup, achieved in part by avoiding recompilation
across a batch, is the same JAX-recompilation cost that
Section~\ref{sec:profiling}'s retained trace analysis reports at the
function level on a different accelerator \citep{openfold2024,colabfold2022}.
AlphaFold2 optimization has thus been a GPU-and-CPU-centric
literature; none of it reports TPU experiments.

Two works come closer to the TPU setting considered here, but address
different questions. ManyFold is a JAX library for
training and validating protein-folding models that does run on TPU,
but the TPU role there is training a different, smaller model
(pLMFold) on a TPU v2-128 pod; its own AlphaFold-style inference-speed
comparison runs on a GPU, not a multi-chip TPU slice
\citep{villegas-morcillo2023manyfold}. MatrixFold describes itself as
the first systematic study of AlphaFold2 mixed-precision inference
\emph{on a many-core ARM CPU architecture}; the framing is close to
this paper's, but the hardware is not
\citep{matrixfold2026}. Neither measures AlphaFold2 inference across a
multi-chip TPU slice using \texttt{pmap}, \texttt{vmap}, or automatic
partitioning, which
is this paper's subject.

JAXBench includes operators extracted from AlphaFold2 among the TPU
kernels it benchmarks \citep{jaxbench2026}; it benchmarks kernels, not
the end-to-end model. A recent characterization of AlphaFold3 profiles
that model's bottlenecks on CPUs and GPUs \citep{iiswc2025af3} and does
not include TPU.

\subsection{Accelerator benchmarking methodology}
\label{sec:rw-benchmarking}

Outside the AlphaFold family, recent trace-based and end-to-end
comparisons of TPU against GPU on language models
\citep{ding2026cclbench,gemma2026tpu} are precedent for comparing the
two through platform-specific stacks on the same model and task; none
touches protein folding, and we do not draw on their results.
Jouppi et al.'s original and TPU v4 architecture papers
\citep{jouppi2017tpu,jouppi2023tpuv4}, together with Google's own v5e
documentation \citep{googletpuv5e}, are our source for the hardware
specifications in Section~\ref{sec:background}, since the v5e
generation used throughout this paper has no dedicated peer-reviewed
architecture paper of its own.

\subsection{Compilation and parallelism frameworks}
\label{sec:rw-frameworks}

GSPMD, the design JAX/XLA's automatic partitioners follow, shards a
computation by propagating sharding decisions outward from explicit
annotations, and its own presentation treats a computation carrying
none as a limiting case of that design rather than as a failure
\citep{xu2021gspmd}. DrJAX documents a
closely related case: removing explicit sharding structure from a
partitioner-based MapReduce system removes its ability to scale efficiently,
even for an embarrassingly parallel pattern \citep{drjax2024}. This is
the closest published precedent to Section~\ref{sec:parallelism}'s
observation, and to our hypothesis that AlphaFold2's unannotated Haiku
modules leave the partitioner nothing to shard on. Neither paper
reports this on AlphaFold2 or on any protein-folding model. We present
our result as consistent with what that design implies when
annotations are absent, not as a new mechanism or a confirmed one. Recent JAX/XLA
measurements on TPU v6e report first-call compilation costs of the
same order as ours \citep{jitcompile2026}, and a separate benchmark
confirms that a new input shape retriggers the cost
\citep{warmup2025}. Finally, Duet Benchmarking and the SPEC
Research Group's methodology for reproducible cloud performance
evaluation document that same-type cloud instances can vary
substantially in measured performance due to hardware heterogeneity
the provider does not expose to the tenant
\citep{bulej2020duet,papadopoulos2021specrg}; we read the gaps
Section~\ref{sec:rh-variability} reports as consistent with that
literature, without a cause isolated, and not as an independent
confirmation of it on ML accelerators.
